\section{Week 3 Refresher: Tokenization and Embedding Geometry}

\subsection*{Setting the Scene}
Before transformers can reason over text, text is mapped to symbols, then to vectors. The mathematics here is mostly linear algebra plus careful distinctions about representation.

\subsection*{Core Reminders}
\begin{itemize}[leftmargin=1.5em]
  \item \textbf{Embedding lookup as linear map:}
  \[
    e_i = E^\top \mathrm{onehot}(i).
  \]
  \item \textbf{BPE/WordPiece/Unigram:} all are segmentation procedures with different optimization heuristics.
  \item \textbf{Cosine similarity:}
  \[
    \cos(u,v)=\frac{u^\top v}{\|u\|\|v\|}.
  \]
  \item \textbf{Norm effects:} same dot product does not imply same cosine ranking.
  \item \textbf{Anisotropy:} embedding spaces can collapse toward dominant directions.
\end{itemize}

\subsection*{Pointer to Earlier Math}
For softmax probabilities and perplexity dependence, see Week 1 refresher.

\subsection*{Common Failure Points}
\begin{itemize}[leftmargin=1.5em]
  \item Mixing sequence-position indices with vocabulary indices.
  \item Treating static embedding analogies as guaranteed behavior in contextual embeddings.
\end{itemize}

\subsection*{Further Refresh}
Jurafsky--Martin (\emph{SLP}), Strang (\emph{Linear Algebra}), CS224N retrieval/embedding lectures.
